% SGA 3, 第 VI_B 篇, 第 1 轮英文重构.
% 范围: 例 5.8.4 全文.
% 底本: Polo--Gille Exp6B-13oct24.pdf, 分篇第 36--37 页
% (合订本第 377--378 页; 重版本页眉 366--367). 第 6 节始于分篇第 37 页后部.
% OCR 与 Reinhold 英文 Markdown 仅为比对材料.
% 本单元无图及底本编者注.

\medskip
\phantomsection\label{ex:sga3-VIB-5.8.4}%
\noindent\textbf{例 5.8.4.---}
沿用前例的记号. 将 \(C\) 描述为序列
\[
  (u_n)_{n\in\mathbb N},
  \qquad u_n\in\mathbb Z/2\mathbb Z,
\]
所成的集合, 可见 \(C\) 去掉一点 \(p\) 后同胚于
\(C\) 的可数个副本的不交并.
事实上, 例如利用 \(C\) 的群结构, 经平移便归结为
\(p=0\) 即零序列的情形. 此时
\[
  C-\{0\}=\coprod_{i\in\mathbb N}C_i,
\]
其中
\[
  C_i=
  \bigl\{(u_n)_{n\in\mathbb N}\mathrel{\big|}
    u_i=1\ \text{且}\ u_j=0\ \text{对}\ j<i\bigr\},
\]
且每个 \(C_i\) 都同胚于 \(C\).
对 \(\lvert F\rvert\) 归纳即可知, 对 \(C\) 的每个非空有限子集
\(F\), 空间 \(C-F\) 同胚于 \(C\) 的可数个副本的不交并,
因而同胚于 \(C\) 去掉一点后的空间.

对每个基数为 \(2\) 的 \(F\subset C\),
设 \(C(F)\) 是 \(C\) 的另一个副本,
以 \(q_F\) 表示 \(C(F)\) 的点 \(0\), 并选取同胚
\[
  \phi_F:C(F)-\{q_F\}\xrightarrow{\sim}C-F.
\]
设 \(X'\) 由将每个 \(C(F)\) 通过 \(\phi_F\) 与 \(C\) 粘合得到.
则 \(X'\) 是不分离的局部 Boolean \(k\)-概形.
此外, 此构造表明, 每个局部常值函数
\[
  f:X'\longrightarrow k
\]
都是常值的. 事实上, 若 \(x,y\in C\) 且 \(F=\{x,y\}\),
则 \(q_F\) 的每个邻域都与 \(x\) 及 \(y\) 的每个邻域相交. 因此
\[
  f(x)=f(q_F)=f(y).
\]
若 \(F'=\{z,t\}\), 其中 \(z,t\in C\), 则同样有
\[
  f(q_{F'})=f(z)=f(q_{\{z,x\}})=f(x).
\]
所以 \(X'\) 连通.

每个点 \(x\in X'\) 也有一个同胚于 \((C,0)\) 的带基点邻域.
更确切地说, 对每个 \(F\), 固定一个带基点拓扑空间的同胚
\[
  h_F:(C,0)\xrightarrow{\sim}(C(F),q_F),
\]
并设 \(T\) 是 \(C\) 的平移群.
若 \(x=q_F\), 使用 \(h_F\); 若 \(x\in C\),
则平移 \(t_x\in T\) 是同胚
\[
  t_x:(C,0)\xrightarrow{\sim}(C,x),
\]
且它是 \(T\) 中有此性质的唯一元素.%
\footnote[1]{\label{note:sga3-VIB-5.8.4-Xprime-hF}译者注:
在这两处, 法文底本分别将刚构造的 \(X'\) 印作 \(X\),
将刚定义的同胚 \(h_F\) 印作 \(q_F\).
英文分别使用 \(X'\) 和 \(h_F\).}

设 \(L\) 是由各 \(h_F\) 生成的自由群, 并置
\[
  G=T*L,
\]
即 \(T\) 与 \(L\) 的自由积 (余积). 对
\[
  h\in\{h_F\}_{\lvert F\rvert=2}\cup T,
\]
用 \(\sigma(h)\) 表示 \(G\) 中相应的生成元,
并以 \(S(h)\) 和 \(B(h)\) 分别表示 \(h\) 的源与目标.
为方便起见, 还令
\[
  \sigma(h_F^{-1})=\sigma(h_F)^{-1},\qquad
  S(h_F^{-1})=B(h_F),\qquad
  B(h_F^{-1})=S(h_F),
\]
并以 \(E\) 表示由 \(T\) 的元素及所有 \(h_F,h_F^{-1}\) 组成的集合.

在乘积空间
\[
  P=G\times X',
\]
上, 其中 \(G\) 赋予离散拓扑, 考虑由
\[
  \bigl(g\sigma(h),x\bigr)\sim\bigl(g,h(x)\bigr)
  \qquad\text{当 }x\in S(h),
\]
对每个 \(h\in E\) 生成的等价关系, 并设 \(Z\) 是商空间.
于是 \(Z\) 由粘合不交并
\[
  \coprod_{g\in G}\{g\}\times X'.
\]
的各个开集而得到. 因此 \(P\) 的每个开集 \(\Omega\)
的饱和化 \(\widetilde\Omega\) 都开
(参见 [BTop], I, \S5.1, 例~2).
明确地说, 由于 \(P\) 的每个开集都是其与各切片
\(\{g\}\times X'\) 之交的并, 只须考虑
\(\{1\}\times W\), 其中 \(W\) 在 \(X'\) 中开.
它的饱和化是 \(\{1\}\times W\) 与集合
\[
  \{\sigma(h_1)\}\times
  h_1^{-1}\bigl(W\cap B(h_1)\bigr),
\]
\[
  \{\sigma(h_1)\sigma(h_2)\}\times
  h_2^{-1}\Bigl(
    h_1^{-1}\bigl(W\cap B(h_1)\bigr)\cap B(h_2)
  \Bigr),
\]
等等的并, 其中遍历 \(E\) 的元素的所有有限序列 \(h_1,\ldots,h_n\).
因此它开, 且投影
\[
  \pi:P\longrightarrow Z
\]
开.

此外, 字
\[
  \sigma(h_1)\cdots\sigma(h_n)
\]
在 \(G\) 中约化, 除非某个 \(h_i\) 是 \(T\) 的单位元,
两个相邻的 \(h_i\) 都属于 \(T\), 或某对 \((h_i,h_{i+1})\) 为
\[
  (h_F,h_F^{-1})
  \quad\text{或}\quad
  (h_F^{-1},h_F).
\]
因此假设 \(x\in X'\), 且
\[
  \beta=
  \bigl(
    \sigma(h_1)\cdots\sigma(h_n),
    h_n^{-1}\cdots h_1^{-1}(x)
  \bigr)
\]
属于 \(\{1\}\times X'\). 此时可假设每个 \(h_i\) 都是一个平移 \(t_i\).
在此情形下,
\[
  \sigma(t_1\cdots t_n)=1
\]
蕴含 \(t_1\cdots t_n=1\), 因而 \(\beta=(1,x)\).
此等价关系与 \(G\) 通过在第一因子上的左平移对 \(P=G\times X'\)
的作用相容. 所以对每个 \(g\in G\),
\(\pi\) 在 \(\{g\}\times X'\) 上的限制单射.

现取 \(z\in Z\), 并选取一个代表元 \((g,x)\in P\).
由前面的讨论,
\[
  U=\pi\bigl(\{g\}\times X'\bigr)
\]
是 \(z\) 的开邻域, 而 \(\pi\) 诱导的连续映射
\[
  \{g\}\times X'\longrightarrow U
\]
开且双射, 因而为同胚.
于是 \(Z\) 局部同构于 \(X'\), 因而也局部同构于 \(C\);
所以它仍是局部 Boolean \(k\)-概形.

最后, \(G\) 在 \(Z\) 上传递地作用.
事实上, 每个 \(z\in Z\) 都与某个形如 \(\pi(1,x)\)
的元素在 \(G\) 下共轭, 因此只须证明每个 \((1,x)\in P\)
都等价于某个 \((\sigma(h),0)\).
若 \(x\in C\), 取 \(h=t_x\); 若 \(x=q_F\), 取 \(h=h_F\).
因此 \(Z\) 是带有离散群 \(G\) 的传递作用的局部 Boolean \(k\)-概形,
但 \(Z\) 不分离.

% 翻译覆盖：本组件全文；源文件 50_expose_VIB_example584_en.tex。
% 自检：行内数学 118，陈列数学 24，原生 tikzcd 0；全部数学及编号保留，公式内自然语言已译。
% 引用键核对：label 2，hyperref 0，SGARefLink 0，行内目标 0；脚注 1，星号脚注正文 0，均与源文件一致。
% 编译及渲染审查由主线程整合完成；此处未宣称编译或独立校对完成。
